\section{Conclusion}
This paper proposes a reciprocal sequential recommendation method, named \emph{ReSeq}. To model dynamic user interests from both sides, we formulate reciprocal recommendation as a distinctive sequence matching task.
Specifically, we encode user behavior sequences from both active and passive perspectives with different module designs.
Subsequently, we conduct comprehensive and fine-grained sequence interactions at both macro and micro levels through multi-scale sequence matching.
In addition, to alleviate the increased time consumption caused by fine-grained interactions, we improve the model deployment efficiency via micro-to-macro self-distillation.
Extensive experiments on real-world datasets show that our approach outperforms several competitive baselines from two perspectives in RRS.
In the future, 
we will consider incorporating multiple auxiliary sequential features (\eg textual information and categories) into our framework in a universal, effective, and efficient way. 
Besides, we would like to further alleviate the sparsity issue while training reciprocal sequential models, including pre-training, data augmentation, and self-supervised learning techniques.

\balance